\section{Introduction}
The Spatial Vessel Analysis project website describes a reproducible workflow for extracting vessel-related features from histological images, with the broader aim of characterizing vascular conformation, remodeling, and spatial relationships to drug distribution. In this context, vessel segmentation is a core preprocessing step because the fidelity of downstream morphometric and spatial analyses depends directly on the quality of the binary vessel mask.

In practical settings, candidate segmentation methods span classical image-processing pipelines, promptable foundation models, and general-purpose neural instance segmentation architectures. These model families differ substantially in their inductive biases, computational cost, and expected transfer performance when applied to histological vascular structures without task-specific retraining.

This repository implements a controlled benchmark of four segmentation approaches on a small image-and-mask dataset stored in TIFF and PNG formats. The compared methods are: (i) a custom VeSpA pipeline implemented as a deterministic color-space and morphology-based segmentation workflow, (ii) the Segment Anything Model (SAM) using automatic mask generation \cite{kirillov2023sam,segmentanything_repo}, (iii) an Ultralytics YOLOv8 segmentation model using the pretrained \texttt{yolov8n-seg.pt} weights \cite{ultralytics_repo}, and (iv) a \textit{novel hybrid VeSpA-SAM approach} that combines the color-based domain knowledge of VeSpA with the segmentation refinement capabilities of SAM using prompt-guided prediction. The hybrid method represents a new strategy for integrating classical domain priors with foundation models in histology.

The objective of this manuscript is not to establish general biological or clinical performance, but to document the repository implementation, standardize evaluation across methods, and report the measured performance obtained from the current codebase. All reported numerical results are taken directly from the benchmark outputs generated by the repository and no additional experiments were introduced beyond those artifacts.
